% ======================================================================
%  Section 4.7.1 --- Nominal case: behaviour without occlusion
%
%  Insert as the FIRST subsection of 4.7, before "Convergence under
%  occlusion" (which becomes 4.7.2).
%
%  Prose that does not depend on the outcome is written as final text.
%  Every point that needs a number from an actual run is marked
%  %%% DATA  and must not be written until the run has been made.
%
%  Requires: the three variants run with  USE_OCCLUDER = false
%            -> results_none.csv, results_tukey.csv, results_hermite.csv
% ======================================================================

\subsection{Nominal case: behaviour without occlusion}
\label{sec:nominal}

Before occlusion is introduced, it is necessary to establish what the three
variants do when there is nothing to reject. Two distinct questions are settled
by this run. The first is whether the weighted formulation is \emph{transparent}
in the absence of outliers: Remark~\ref{rem:baseline} asserts that
$\mathbf{D}=\mathbf{I}$ reduces~\eqref{eq:robust-lm-law} to the control law of
Chapter~\ref{ch:herpvs}, and the nominal run verifies numerically what that
remark establishes algebraically. The second is a matter of attribution: unless
the three variants are known to agree when the scene is clean, any difference
observed later under occlusion cannot be ascribed to the treatment of the
occlusion rather than to the weighting disturbing the nominal behaviour of the
control law.

\subsubsection*{Conditions}

The reference image $\mathbf{I}^{*}$ is acquired from the unoccluded target at
the desired pose, as in every configuration of this chapter. For the nominal run
alone, the images acquired during servoing are also taken from the unoccluded
target, so that the residual arises solely from the pose error and no
measurement is corrupted. The initial pose is
%%% DATA: state the initial pose actually used, e.g.
%%% $(t_x,t_y,t_z)=(0.02,-0.02,1.36)\,\mathrm{m}$,
%%% $(\theta u_x,\theta u_y,\theta u_z)=(2,-2,1)^{\circ}$,
%%% and the corresponding mean image displacement with respect to the
%%% desired view (in pixels), which fixes the difficulty of the task.
and the desired pose is the frontal view at $Z=1.3\,\mathrm{m}$ at which the
target fills the image. All remaining quantities --- the Hermite bank and its
four scales, the interaction matrix~\eqref{eq:stacked-L}, the adaptive gain and
damping~\eqref{eq:lambda}--\eqref{eq:mu}, the Tukey constant, the sampling
period and the stopping criterion --- are those fixed in
Section~\ref{sec:ablation-protocol} and are identical across the three variants.

\subsubsection*{Reported quantities}

Figure~\ref{fig:nominal-error} superimposes the evolution of the feature-error
norm $\norm{\mathbf{e}}$ for the three variants, and the corresponding final
pose errors and iteration counts are given in the first block of
Table~\ref{tab:ch4-ablation}. Transparency of the weighting is taken to hold
when the three trajectories are indistinguishable to within the numerical
tolerance of the solver, and when the final pose errors agree to the precision
reported.

%%% DATA: the outcome. Two cases; write whichever the runs show, and do not
%%% write either before they have been made.
%%%
%%% CASE A --- the three trajectories coincide:
%%%   "The three trajectories coincide throughout, and the final pose errors
%%%    agree to <precision>. The weighting is therefore transparent in the
%%%    absence of outliers: with no residual large enough to be rejected, the
%%%    Tukey weights remain at unity across the image and both weighted
%%%    variants reproduce the unweighted control law. Remark~\ref{rem:baseline}
%%%    is thus confirmed numerically, and any difference observed in the
%%%    occluded configurations of Section~\ref{sec:ablation-results} may be
%%%    attributed to the treatment of the occlusion alone."
%%%
%%% CASE B --- Variant 2 (and possibly 3) departs from Variant 1:
%%%   This is not a defect and must not be treated as one. It is the
%%%   phenomenon anticipated in Section~\ref{sec:tukey-limitation}: while the
%%%   initial pose error is large, valid measurements in photometrically
%%%   sensitive regions produce legitimately large residuals, and a
%%%   redescending estimator rejects a proportion of them. The nominal run then
%%%   ceases to be a mere verification and becomes a first result, namely that
%%%   residual-based rejection carries a cost even when the scene is clean.
%%%   Report the magnitude of that cost (difference in iterations and in final
%%%   pose error), and state whether the structure normalization of
%%%   Section~\ref{sec:hermite-weight} reduces it --- which is precisely what
%%%   it was designed to do, since it shrinks the residual where the local
%%%   Hermite energy is high.

\begin{figure}[t]
    \centering
    % \includegraphics[width=\linewidth]{figures/ch4_nominal_error.pdf}
    \caption{Nominal case, unoccluded target: evolution of the feature-error
    norm $\norm{\mathbf{e}}$ against iteration for the three weighting variants
    of Section~\ref{sec:ablation-protocol}. With no corrupted measurement
    present, the weighted variants are expected to reproduce the unweighted
    control law.}
    \label{fig:nominal-error}
\end{figure}
